\scanpage{158}
\begin{center}- 111 -\end{center}\hfill IX

Reste à décrire des isomorphismes d). En fait, sur une base quelconque $U$, on a un homomorphisme canonique
\[
 (9.4.3.1)\qquad \mathbb G_m(U)\otimes\check H_{\unsure{o|0}}\longrightarrow H^1(U,D(H_{\unsure{o|0}U}))\ ,
\]
qu'on décrit, lorsque $H_{\unsure{o|0}}$ est annulé par l'entier $m\geq0$, à l'aide de l'homomorphisme cobord $\mathbb G_m(U)\xrightarrow{\partial}H^1(U,\mu_m)$ associé à la suite exacte de Kummer, en prenant le composé
\diagram{\begin{tikzcd}[column sep=small]
 \mathbb G_m(U)\otimes\check H_{\unsure{o|0}}\arrow[r]& H^1(U,\mu_m)\otimes\check H^{\unsure{o|0}}\arrow[r]& H^1(U;\mu_m\otimes\check H_{\unsure{o|0}})\arrow[r]&H^1(U,D(H_{\unsure{o|0}S}))\ ,
\end{tikzcd}}
où la deuxième flèche se déduit de l'homomorphisme évident $\check H_{\unsure{o|0}}\longrightarrow\operatorname{Hom}(\mu_m,\mu_m\otimes\check H_{\unsure{o|0}})$, et la troisième est un cas particulier de b) (où on remplace $n$ par $m$). Nous laissons au lecteur le soin de vérifier que (9.4.3.1) ne dépend pas du choix de l'entier $m$ tel que $mH_{\unsure{o|0}}=0$.
D'autre part je dis que l'homomorphisme (9.4.3.1) est toujours injectif, et que son conoyau est nul si $\operatorname{Pic}(U)=0$. (En fait, le conoyau est toujours isomorphe à $\operatorname{Hom}(H_{\unsure{o|0}},\operatorname{Pic}(U))$, mais peu importe ici.) Pour voir ceci, on est ramené en effet au cas où $H_{\unsure{o|0}}$ est de la forme $\mathbb Z/n\mathbb Z$, et alors l'assertion résulte aussitôt de la suite exacte de cohomologie associée à la suite exacte de Kummer (*) ci-dessus. Ceci prouve la formule d). On en conclut aussitôt :

\uline{Corollaire} 9.4.4. \uline{Avec les notations de} 9.4.3 \uline{on a un isomorphisme canonique}
\[
 \text{e)}\qquad H^1(\eta,D(H)_\eta)/\operatorname{Im}H^1(S,D(H)\simeq\check H_{\unsure{o|0}}\ ,
\]
\uline{d'où} (\uline{en combinant avec} a) \uline{et} b)) \uline{un isomorphisme canonique}
\figure{\begin{tikzpicture}\draw (0,0)--(.15,.25)--(.35,0);\end{tikzpicture}}
\[
 (9.4.5)\qquad\operatorname{Ext}^1(Q_\eta,P_\eta)/\operatorname{Im}\operatorname{Ext}^1(Q,P)\simeq\check Q_{\unsure{o|0}}\otimes\check Q_{\unsure{o|0}}^*\ .
\]
\margin{\figure{\begin{tikzpicture}\draw (0,.1)--(.45,.1) (.35,.28)--(.38,0);\draw (.85,.18)--(.9,0);\draw (5.5,0)--(5.6,.17)--(5.7,0);\end{tikzpicture}}}
\margin{429}
\margin{429}
\margin{157}
